\documentclass[conference]{IEEEtran}
\IEEEoverridecommandlockouts
% The preceding line is only needed to identify funding in the first footnote. If that is unneeded, please comment it out.
%Template version as of 6/27/2024

\usepackage{amsmath,amssymb,amsfonts}
\usepackage{algorithmic}
\usepackage{graphicx}
\usepackage{textcomp}
\usepackage{xcolor}
\usepackage[backend=biber,style=numeric]{biblatex}
\addbibresource{assets/latex/report.bib} 

 \graphicspath{{assets/latex/}}

\begin{document}

\title{CS6524 Computer Vision Lab. 2\\ Camera Calibration and Stereo Vision}

\author{\IEEEauthorblockN{1\textsuperscript{st} Dawid Jakubowski}
\IEEEauthorblockA{\textit{Immersive Software Engineering} \\
\textit{University of Limerick}\\
Limerick, Ireland \\
23384549@studentmail.ul.ie}
\and
\IEEEauthorblockN{2\textsuperscript{nd} Given Name Surname}
\IEEEauthorblockA{\textit{dept. name of organization (of Aff.)} \\
\textit{name of organization (of Aff.)}\\
City, Country \\
email address or ORCID}
\and
\IEEEauthorblockN{2\textsuperscript{nd} Given Name Surname}
\IEEEauthorblockA{\textit{dept. name of organization (of Aff.)} \\
\textit{name of organization (of Aff.)}\\
City, Country \\
email address or ORCID}
}

\maketitle

\section{Introduction}
This lab investigates camera calibration and stereo vision for estimating depth from video recordings. Camera intrinsic parameters and lens distortion coefficients are first determined using checkerboard images. These parameters are then used to estimate the relative camera motion between consecutive video frames. Finally, stereo rectification and disparity estimation are used to reconstruct scene depth.

\section{Methodology}
\subsection{Step 1}
\subsection{Step 2}
\subsection{Step 3}
\subsection{Step 4: Recording scene video}
Two videos were recorded using the calibrated smartphone camera, capturing static scenes containing objects at different depths. The camera was moved horizontally to introduce a baseline between consecutive frames while minimising rotation. The recordings were subsequently resized and aligned with the calibration image dimensions to ensure consistency during processing.


\subsection{Step 5: Inter-Frame Translation & Pose Matrices}
SIFT was used to detect and describe keypoints in consecutive video frames. Corresponding features were matched using a brute-force matcher and Lowe's ratio test with a threshold of 0.75. After correcting for lens distortion, the Essential Matrix was estimated using RANSAC and decomposed using cv2.recoverPose() to obtain the relative rotation matrix and translation direction. The Fundamental Matrix was derived using the camera intrinsics. Unreliable estimates were flagged based on pose inlier counts, inlier ratios, and rotation thresholds.

\subsection{Step 6}


\section{Discussion and Failure Analysis}

\subsection{Citations}

If you need to cite external sources, use the \texttt{assets/latex/report.bib} file. As an example, here's how you cite the textbook~\cite{foundationsCVbook}.

\printbibliography

\end{document}
